\documentclass{resume} % Use the custom resume.cls style
\usepackage[none]{hyphenat}
\usepackage{xcolor}
\usepackage{hyperref}
\usepackage{graphicx}
\hypersetup{
  colorlinks=true,
  linkcolor=blue,
  filecolor=magenta,      
  urlcolor=blue,
  pdftitle={Cristian Barrera Resume},
  pdfpagemode=UseOutlines,
  }

\newcommand{\tab}[1]{\hspace{.2667\textwidth}\rlap{#1}}
\newcommand{\itab}[1]{\hspace{0em}\rlap{#1}}

\name{Cristian Barrera}
\email{cris.rbarreram@gmail.com}
\phone{+1 (216) 333-4388}
\website{crisbarrera.com}
\companywebsite{bartekllc.org}
\address{Cleveland, Ohio}
\linkedin{linkedin.com/in/cristian-barrera-63205997/}
\github{github.com/maberyick}

\begin{document}
%%%
%%%
%%%
\begin{rSection}{SUMMARY}
Data Scientist focused on medical imaging and machine learning, with 7+ years across academic research, biotech, and production clinical pipelines. Strong experience in OCT, deep learning, distributed training, and MLOps, with a track record of shipping practical tools used by research and clinical teams.
\end{rSection}
%%%
%%%
%%%
\begin{rSection}{EDUCATION}
\begin{rSchool}{Georgia Tech \& Emory University, Atlanta, GA}
  {\textbf{Aug 2020} -- \textbf{Jan 2026}}
  {Ph.D., Biomedical Engineering}
  {Medical Imaging, Histopathology, Ophthalmology}
  {}
\end{rSchool}
%
\begin{rSchool}{National University of Colombia}
  {\textbf{Jan 2017} -- \textbf{Dec 2019}}
  {Master of Engineering, Biomedical Engineering}
  {Medical Image Processing, Machine Learning, Genomics}
  {}
\end{rSchool}
%
\begin{rSchool}{University of Wales, London, UK}
  {\textbf{2015} -- \textbf{2016}}
  {Diploma, Business IT}
  {}
  {}
\end{rSchool}
%
\begin{rSchool}{South Colombian University}
  {\textbf{Jan 2011} -- \textbf{Dec 2015}}
  {Bachelor of Engineering, Electronics}
  {Digital Electronics, Signal Processing, Telecommunications}
  {}
\end{rSchool}
\end{rSection}
%%%
%%%
%%%
\begin{rSection}{EXPERIENCE}

  \begin{rWork}{Data Scientist III}{Cleveland Clinic (Cole Eye Institute) - Cleveland, OH}{Nov 2024 -- Present}
    \item Lead OCT foundation model development by pretraining iVisionFM on 3.2M B-scans, then fine-tuning on hundreds of thousands to millions of scans for layer-specific tasks (ILM, BM, RPE, and related boundaries).
    \item Designed a mixture-of-experts architecture with a pathology-aware router, delivering 30--40\% performance improvement versus a TransUNet-only baseline on complex pathology cohorts.
    \item Architected distributed training and inference across 20+ GPUs on 10+ Linux workstations with containerized execution.
    \item Built an Airflow-based orchestration workflow for OCT run submission across multiple computers, integrating 20+ containerized models.
    \item Delivered a React and PostgreSQL platform for run submission, tracking, and quality-control visibility at 50--100 runs per day.
  \end{rWork}

  \begin{rWork}{Owner}{BARTEK LLC (bartekllc.org) - Cleveland, OH}{2024 -- Present}
    \item Founded BARTEK LLC and delivered a 9-month Veterans Affairs contract (Sep 2025 -- Jun 2026) to build a histopathology image processing platform.
    \item Built a web workflow for users to submit histopathology runs, monitor progress, and review outputs.
    \item Developed a containerized multi-model pipeline with foundational models and run-allocation logic across compute resources.
    \item Implemented an AI assistant to recommend the best models to execute for each submission context.
  \end{rWork}
  
  \begin{rWork}{Graduate Research Scientist}{Emory University - Atlanta, GA}{Feb 2018 -- Jan 2024}
    \item Designed data models and visualizations for 4,000+ histopathology images for cancer and biomarker research.
    \item Developed machine learning models for patient outcomes and immunotherapy response contributing to peer-reviewed publications.
    \item Built scalable deep learning and traditional ML pipelines for biomedical imaging and clinical datasets.
    \item Collaborated with multidisciplinary teams to translate computational results into clinically relevant findings.
  \end{rWork}

  \begin{rWork}{Imaging Scientist Intern}{Genentech - South San Francisco, CA}{Jun 2023 -- Aug 2023}
    \item Improved data processing efficiency by 30\% through Python-based workflows.
    \item Developed ETL patterns for immunofluorescence imaging to integrate heterogeneous research datasets.
    \item Partnered with scientists to deliver reproducible analysis outputs for imaging studies.
  \end{rWork}

  \begin{rWork}{Visiting Scientist}{Case Western Reserve University - Cleveland, OH}{Nov 2017 -- Jan 2020}
    \item Supported imaging and data analysis workflows in biomedical engineering research collaborations.
    \item Applied deep learning and signal processing methods to experimental datasets.
    \item Delivered reproducible code and documentation for faculty and lab teams.
  \end{rWork}

  \begin{rWork}{Biomedical Engineering Contractor}{National University of Colombia - Bogota, Colombia}{Jan 2017 -- Dec 2019}
    \item Developed deep learning face recognition tools for longitudinal child tracking programs.
    \item Designed data collection and evaluation workflows to monitor accuracy and false alerts in deployment.
    \item Supported surface EMG testing for maternal-fetal heart rate capture and downstream signal analysis.
  \end{rWork}

\end{rSection}
%%%
%%%
%%%
\begin{rSection}{PROJECTS \& CURRENT WORK}
\textbf{OCT Foundation Model and MoE Segmentation Stack.} Pretrained iVisionFM on 3.2M OCT B-scans and fine-tuned for retinal layer segmentation (including ILM, BM, and RPE) with hundreds of thousands to millions of labeled examples, then added a pathology-aware mixture-of-experts router to improve complex-case performance by 30--40\% versus a TransUNet-only approach.

\vspace{2mm}

\textbf{Ophthalmology OCT Workflow Platform.} Distributed Airflow and container-based workflow for OCT processing across multiple computers, integrating 20+ models behind a React interface for run submission, monitoring, and quality control, with PostgreSQL tracking 50--100 daily runs.
\end{rSection}

%\begin{rSection}{PROJECTS}
%\textbf{Hyperplex Tissue Imaging Enhancement.} Led development of Python-based image analysis workflows for multiplexed microscopy, improving accuracy and efficiency by 25\%. Enhanced HPC deployment processing speed by 20\%.

%\vspace{2mm}

%\textbf{PhenopyCell: Lung Cancer Prediction.} Developed AI-powered system analyzing large pathology datasets, improving diagnostic accuracy by 20\%. Applied CNNs to reduce errors by 15\% and improve speed by 25\%.

%\end{rSection}
%%%
%%%
%%%
\begin{rSection}{SKILLS}
\textbf{Main Languages:} Python, R

\textbf{Frameworks/Tools:} PyTorch, TensorFlow, Airflow, React, Flask, Docker, Podman, Redis, PostgreSQL, MLflow, Git, Linux, HPC, AWS, Azure

\textbf{Domains:} Medical Imaging, OCT, Histopathology, Machine Learning, Deep Learning, Data Engineering, MLOps, Ophthalmology

\end{rSection}
%%%
%%%
%%%
\end{document}
